\chapter{Manual de instalación y despliegue}
\label{app:manual_instalacion}

\section{Requisitos previos}
\label{sec:requisitos_previos}

\begin{itemize}
    \item \textbf{Sistema operativo:} el entorno utilizado durante el
    desarrollo y validación del sistema ha sido Ubuntu 24.04 LTS. También es
    posible utilizar otros sistemas compatibles con Docker.

    \item \textbf{Lenguaje y runtime:} no es necesario instalar manualmente
    Python ni sus dependencias en el sistema anfitrión, ya que el entorno de
    ejecución se encuentra definido dentro de los contenedores del proyecto.

    \item \textbf{Gestor de paquetes:} las dependencias necesarias para
    ejecutar la aplicación y el generador se instalan automáticamente durante
    la construcción de las imágenes Docker.

    \item \textbf{Servicios externos:} no es necesario instalar manualmente
    servicios adicionales. La configuración Docker del proyecto incluye los
    contenedores necesarios para los servicios utilizados por UVLHub, entre
    ellos la base de datos MariaDB, Redis, Elasticsearch, MailHog y Selenium.

    \item \textbf{Otras herramientas:} es necesario disponer de Git para
    obtener el código fuente y Docker junto con Docker Compose para construir y
    ejecutar el entorno de desarrollo.
\end{itemize}

\section{Disponibilidad del código}
\label{sec:disponibilidad_codigo}

El código fuente relacionado con este trabajo está disponible públicamente en
GitHub. La Tabla~\ref{tab:repositorios_codigo} recoge la dirección de cada
repositorio, el commit de referencia considerado en esta memoria y la licencia
declarada en cada caso.

\begin{table}[H]
\centering
\small
\setlength{\tabcolsep}{4pt}
\renewcommand{\arraystretch}{1.2}
\begin{tabularx}{\textwidth}{@{}
>{\raggedright\arraybackslash}p{2.5cm}
>{\raggedright\arraybackslash}X
>{\raggedright\arraybackslash}p{5.2cm}
>{\centering\arraybackslash}p{2.1cm}
@{}}
\toprule
\textbf{Repositorio} & \textbf{URL} & \textbf{Commit de referencia} &
\textbf{Licencia} \\
\midrule

\texttt{uvlhub}
&
\url{https://github.com/antjimort/uvlhub}
&
{\footnotesize\ttfamily\shortstack[l]{b4ca8d81f5b03c688a14\\68289a74d144ce5d081a}}
&
MIT
\\

\texttt{uvlhub\_v2}
&
\url{https://github.com/antjimort/uvlhub\_v2}
&
{\footnotesize\ttfamily\shortstack[l]{45ee7a8c8f0ed0b800ee\\7a4646a561ae0145995f}}
&
No declarada
\\

\texttt{fm\_generator}
&
\url{https://github.com/antjimort/fm\_generator}
&
{\footnotesize\ttfamily\shortstack[l]{5b0a32f7d728defcfcaf\\736e6c1f88d3d5e3c171}}
&
No declarada
\\

\bottomrule
\end{tabularx}
\caption{Repositorios, commits de referencia y licencias declaradas.}
\label{tab:repositorios_codigo}
\end{table}

El repositorio \texttt{uvlhub} contiene la aplicación web y su integración con
el generador. En su directorio raíz se ha añadido un archivo \texttt{LICENSE}
que declara la licencia MIT. Se mantienen los avisos de licencia específicos
incluidos en determinados archivos del proyecto. La licencia del repositorio
\texttt{uvlhub} no se extiende automáticamente a los repositorios
\texttt{uvlhub\_v2} y \texttt{fm\_generator}, que no tienen una licencia
declarada.

La aplicación está desplegada en Render y puede consultarse en
\url{https://uvlhub-antjimort-8e4j.onrender.com}. La versión disponible es la
\texttt{v1.0.0}, cuyo código ejecutable corresponde al commit
\texttt{b4ca8d81f5b03c688a1468289a74d144ce5d081a} de \texttt{uvlhub}. El
commit posterior que añade el archivo de licencia no modifica la aplicación.
El generador se ejecuta como parte de UVLHub y no dispone de un despliegue
público independiente. Los repositorios \texttt{uvlhub\_v2} y
\texttt{fm\_generator} tampoco cuentan con despliegues separados.

Además del código disponible en los repositorios indicados, se ha depositado
en Zenodo una copia archivada de \textit{UVLHub with fm\_generator integration}.
El archivo corresponde a UVLHub 1.0.0 e incluye la versión 0.0.1 de
\texttt{fm\_generator}. Esta versión puede citarse mediante el DOI
\href{https://doi.org/10.5281/zenodo.22963634}{10.5281/zenodo.22963634};
su registro está disponible en
\url{https://zenodo.org/records/22963634}.

\section{Instalación de dependencias y ejecución}
\label{sec:instalacion_dependencias}


Las dependencias del proyecto se gestionan mediante Docker, por lo que no es
necesaria una instalación manual de paquetes. La construcción de los
contenedores instala automáticamente las dependencias definidas en la
configuración del proyecto.

El siguiente comando construye las imágenes necesarias y levanta el entorno de
desarrollo definido para UVLHub. Durante este proceso se instala también el
paquete \texttt{fm\_generator} y se construyen los recursos necesarios para su
ejecución desde el navegador mediante Pyodide.

\begin{verbatim}
docker compose -f docker/docker-compose.dev.yml up --build
\end{verbatim}

Para detener y eliminar los contenedores se usa el siguiente comando:

\begin{verbatim}
docker compose -f docker/docker-compose.dev.yml down
\end{verbatim}

\section{Despliegue en producción}
\label{sec:despliegue_produccion}


La herramienta utilizada para el despliegue de la app principal ha sido Render, pero ha sido usada únicamente para las revisiones de los tutores. También se ha usado Redis, utilizado como almacenamiento externo para la gestión de las
sesiones de Flask para conservar el estado
del proceso de configuración del generador.

No obstante, la plataforma utilizada no constituye un requisito del sistema,
ya que la aplicación puede ser desplegada mediante otras alternativas
compatibles con contenedores.

El despliegue oficial de UVLHub es gestionado por DiversoLab.


\section{Verificación de la instalación}
\label{sec:verificacion-instalacion}


Una vez levantados los contenedores mediante el comando indicado en la
Sección~\ref{sec:instalacion_dependencias}, se puede comprobar que la
instalación se ha realizado correctamente accediendo a la dirección local
proporcionada por la aplicación durante el arranque.

La dirección de acceso local depende de los puertos publicados por la
configuración Docker. Puede comprobarse mediante el siguiente comando:

\begin{verbatim}
docker compose -f docker/docker-compose.dev.yml ps
\end{verbatim}

El servicio debe abrirse en el navegador utilizando el puerto expuesto en el
equipo anfitrión, normalmente mediante una dirección de la forma
\texttt{http://localhost:<PUERTO>}. La salida del arranque permite comprobar
que el contenedor web se ha iniciado correctamente; véase la Figura~\ref{fig:inicio-contenedor}.

\begin{figure}[H]
    \centering
    \includegraphics[width=0.9\textwidth]{salida_contenedor.png}
    \caption{Salida de consola tras el arranque correcto del contenedor web. Elaboración propia.}
    \label{fig:inicio-contenedor}
\end{figure}

\subsection{Ejecución de las pruebas automatizadas}

Una vez iniciado el entorno Docker, la suite de pruebas integrada del módulo
\textit{generator} puede ejecutarse mediante el siguiente comando:

\begin{verbatim}
docker compose -f docker/docker-compose.dev.yml exec web rosemary test generator --all
\end{verbatim}

La ejecución correcta debe finalizar sin errores. Esta suite valida el flujo del
asistente, la generación de modelos UVL, la persistencia de parámetros y las
opciones de salida integradas en UVLHub.

Las pruebas del paquete independiente \texttt{fm\_generator} se ejecutan desde
la raíz del repositorio mediante:

\begin{verbatim}
pytest fm_generator/tests -v
\end{verbatim}

\subsection{Medición de cobertura}

La cobertura del código del módulo \textit{generator} integrado en UVLHub puede
obtenerse ejecutando la suite mediante Rosemary dentro del contenedor web:

\begin{verbatim}
docker compose -f docker/docker-compose.dev.yml exec web rosemary coverage generator --all 
\end{verbatim}

El informe generado incluye tanto el código de producción como los ficheros de
prueba. Para los resultados recogidos en esta memoria se considera únicamente
la cobertura de los ficheros de producción, tal y como se explica en la
Sección~\ref{sec:cobertura}.

La cobertura del paquete independiente \texttt{fm\_generator} puede obtenerse
mediante el siguiente comando:

\begin{verbatim}
pytest fm_generator/tests -v --cov=fm_generator --cov-report=term-missing --cov-report=html
\end{verbatim}

El informe textual muestra las líneas no cubiertas, mientras que el informe HTML
se genera en el directorio \texttt{htmlcov}. Los resultados de cobertura
utilizados en esta memoria se describen en la
Sección~\ref{sec:cobertura}.

